% Generated draft: requires mathematical and source review.
\documentclass[11pt]{article}
\usepackage[margin=25mm]{geometry}
\usepackage{fontspec}
\setmainfont{Latin Modern Roman}
\usepackage{amsmath,amssymb,mathtools,unicode-math}
\setmathfont{Latin Modern Math}
\usepackage{enumitem,needspace,xurl,hyperref}
\setlength{\parindent}{0pt}
\setlength{\parskip}{5pt}
\setlength{\emergencystretch}{3em}
\newcommand{\N}{\mathbb{N}}
\newcommand{\Z}{\mathbb{Z}}
\newcommand{\Q}{\mathbb{Q}}
\newcommand{\R}{\mathbb{R}}
\newcommand{\C}{\mathbb{C}}
\newcommand{\F}{\mathbb{F}}
\newcommand{\A}{\mathbb{A}}
\newcommand{\PP}{\mathbb{P}}
\newcommand{\E}{\mathbb{E}}
\newcommand{\eps}{\varepsilon}
\newcommand{\dd}{\,\mathrm{d}}
\begin{document}
\section*{UnsolvedMath ID: 30002577}
\textbf{Source problem code:} catalogue entry
\subsection{Definitions and mathematical statement}
Let \(k\) denote the ground field. The supplied record does not specify its characteristic, whether it is algebraically closed, or whether the questions are intended uniformly over all fields; these remain unresolved parameters. All tensor products below are over \(k\). A Hopf algebra is a unital associative \(k\)-algebra equipped with a compatible comultiplication, counit, and antipode. For this draft, ``noetherian'' means both left and right noetherian; the intended sidedness must be checked against the source.

For an algebra \(A\), let \(\operatorname{MinSpec}(A)\) denote the set of two-sided prime ideals minimal under inclusion. Here a proper two-sided ideal \(P\) is prime if \(IJ\subseteq P\) implies \(I\subseteq P\) or \(J\subseteq P\) for two-sided ideals \(I,J\). Two unital algebras are Morita equivalent if their categories of unital left modules are equivalent. The record does not specify whether this equivalence must be \(k\)-linear. An algebra is PI if it satisfies a nonzero polynomial identity with coefficients in \(k\). For a nonzero unital algebra, its Gelfand--Kirillov dimension is
\[
\operatorname{GKdim}_k A
=\sup_V\limsup_{n\to\infty}
\frac{\log\dim_k(V^n)}{\log n},
\]
where \(V\) ranges over finite-dimensional \(k\)-subspaces of \(A\) containing \(1\), and \(V^n\) is the span of products of \(n\) elements of \(V\). This definition allows infinite dimension and does not assume that \(A\) is finitely generated. Call \(A\) strongly noetherian if \(A\otimes_k C\) is left and right noetherian for every commutative noetherian unital \(k\)-algebra \(C\); the source's precise convention also requires confirmation.

The supplied record combines the following questions.
\begin{enumerate}
\item For every noetherian Hopf algebra \(H\), is it true that
\[
\forall P,Q\in\operatorname{MinSpec}(H),\qquad
H/P\text{ and }H/Q\text{ are Morita equivalent}?
\]
If this fails, determine which properties the minimal-prime quotients nevertheless share. In particular, determine whether, for every such \(H\) and every \(P,Q\in\operatorname{MinSpec}(H)\),
\[
H/P\text{ is PI}\ \Longleftrightarrow\ H/Q\text{ is PI},
\qquad
\operatorname{GKdim}_k(H/P)=\operatorname{GKdim}_k(H/Q).
\]
These are proposed comparison questions, not asserted equalities. The quotients are considered as algebras: the record does not assert that the minimal prime ideals are Hopf ideals.

\item Is every noetherian Hopf algebra strongly noetherian? With the conventions above, the requested assertion is
\[
\forall H\;\bigl[H\text{ is a noetherian Hopf }k\text{-algebra}
\Longrightarrow
\forall C\;\bigl(C\text{ is a commutative noetherian }k\text{-algebra}
\Longrightarrow H\otimes_k C\text{ is noetherian}\bigr)\bigr].
\]
The range of permissible ground fields remains unspecified.

\item Construct and classify Hopf-algebra extensions of enveloping algebras by finite-dimensional pointed Hopf algebras, with fixed \(H\) or fixed Gelfand--Kirillov dimension, and relate or describe their root systems. Here \(U(\mathfrak g)\) denotes the universal enveloping algebra of a Lie algebra \(\mathfrak g\), and a Hopf algebra is pointed when every simple subcoalgebra is one-dimensional. The record does not identify the role of \(H\) in this clause, the direction or exactness convention for the extension, the allowed Lie algebras and pointed Hopf algebras, the equivalence relation for classification, the permitted dimension values, or the intended root systems. Consequently this clause specifies a research task but does not determine a unique classification problem.

\item Can every known finite-dimensional Nichols algebra be constructed from a diagonal Nichols algebra by diagram folding, and do folded liftings give all liftings? A braided vector space is a vector space \(V\) with an invertible linear map \(c:V\otimes V\to V\otimes V\) satisfying the braid equation. Its Nichols algebra \(\mathcal B(V,c)\) is the quotient of the braided tensor Hopf algebra \(T(V)\) by its largest homogeneous braided Hopf ideal contained in degrees at least two. The braiding is of diagonal type if some basis \((x_i)\) satisfies
\[
c(x_i\otimes x_j)=q_{ij}x_j\otimes x_i,
\qquad q_{ij}\in k^\times.
\]
The qualifier ``known'' is retained: no collection of examples or knowledge cutoff is specified. The record does not define diagram folding, the ambient braided category, the admissible folding data, the associated graded object defining a lifting, or equivalence of liftings. Thus neither the construction claim nor the exhaustion claim has a fully specified quantifier domain. No particular folding or lifting convention is selected here.
\end{enumerate}
\subsection{Short English statement}
The record asks whether minimal-prime quotients of a noetherian Hopf algebra are Morita equivalent or share PI status and Gelfand--Kirillov dimension, and whether every noetherian Hopf algebra is strongly noetherian. It also asks for classifications of certain enveloping-algebra extensions and for folding constructions of known finite-dimensional Nichols algebras and their liftings. Ground-field assumptions and several essential classification, folding, and lifting conventions are missing. The supplied open status is dated curation, and this statement-only draft requires human mathematical review.
\subsection{Sources}
\begin{itemize}
\item Supplied local UnsolvedMath registry record, permanent ID \(30002577\), exact source code \texttt{catalogue entry}, title ``Structure and Classification Problems for Noetherian Hopf Algebras,'' in the collection ``Oberwolfach Reports Open Problems.'' No individual proposer is supplied. The proposed year is recorded as 2014.
\item Supplied report citation: \emph{Mini-Workshop: Infinite Dimensional Hopf Algebras (2015)}, Oberwolfach Reports, pp.~1111--1137, DOI \texttt{10.4171/owr/2014/20}. \url{https://doi.org/10.4171/owr/2014/20}. The background attributes the questions to workshop open-problem discussions. The citation's 2015 date and the record's 2014 proposed year require reconciliation.
\item Supplied literature source, labelled ``Oberwolfach open-questions report'': \url{https://doi.org/10.4171/owr/2014/20}.
\item Supplied literature source, labelled ``Finite-GK pre-Nichols algebras'': \url{https://arxiv.org/abs/2504.12206}.
\item Supplied literature source, labelled ``Pointed Hopf algebras and finiteness'': \url{https://arxiv.org/abs/2507.23730}.
\item Literature assessment dated exactly \texttt{2026-08-21}: the supplied status is ``open.'' The supplied assessment says that important diagonal and low-growth cases have been classified, while the combined questions concerning minimal-prime quotients, strong noetherianity, non-diagonal extensions, diagram folding, and all liftings retain unresolved general cases. This is supplied dated curation derived from the OWR record, not an independently verified literature assessment. The background expressly requests independent verification of time-sensitive status and cited claims.
\item The registry labels the statement \texttt{corrected\_verified} and says that notation, typography, and wording were normalized without changing the mathematical content. These are supplied editorial claims, not independent mathematical verification. No report-internal problem location or individual attribution is supplied.
\end{itemize}
\end{document}
